% Source : cours/cm8.qmd. Effets conditionnels de la consolidation.
\begin{board}{8}{Austérité : le ratio dette / PIB peut augmenter}
 {Le résultat dépend du multiplicateur, de la dette initiale et du contexte.}
\stage{5.4}{1}{La dépense est réduite}
\stage{16}{2}{Des effets en chaîne}
\stage{26.6}{3}{Le dénominateur diminue}
\bill{5}{8.5}
\draw[drawline] (2.7,11.5) circle (.4) (2.7,10.5) circle (.4);
\draw[drawline] (3.05,11.25)--(6.6,9.6) (3.05,10.75)--(6.6,12.1);
\node[label,text width=3cm] at (7.8,8.8) {$G\ \downarrow$};
\captext{5.4}{Moins d'emprunt, à autres\\effets inchangés.}
\factory{15.9}{9}
\person{13}{6.6}{blue}
\shop{18.4}{6.6}
\draw[flow] (14.6,7.7)--(15.9,7.7);
\captext{16}{Revenus, consommation\\et PIB diminuent.}
\node[text=ink,font=\fontsize{43}{50}\selectfont,inner sep=0pt] at (26.6,10.5)
 {$\dfrac{\textnormal{Dette}}{\textnormal{PIB}\ \downarrow}$};
\draw[possible] (29.7,7.2)--(29.7,9);
\noteat{26.6}{6.4}{Le ratio peut monter}
\captext{26.6}{Si le PIB baisse assez,\\le ratio peut augmenter.}
\end{board}
